\section*{A}

\textbf{A/B Testing:} A method of comparing two versions of a model or system to determine which performs better.

\textbf{Adam:} Adaptive Moment Estimation, a popular optimization algorithm for training neural networks.

\textbf{Adversarial Attack:} An attempt to fool an ML model by providing deceptive input.

\textbf{API (Application Programming Interface):} A set of protocols for building and integrating application software.

\textbf{AutoML:} Automated Machine Learning, tools and techniques for automating the ML workflow.

\section*{B}

\textbf{Backpropagation:} Algorithm for training neural networks by computing gradients through the chain rule.

\textbf{Batch Inference:} Making predictions on large batches of data offline.

\textbf{Batch Normalization:} A technique to normalize layer inputs to stabilize neural network training.

\textbf{Bias (Statistical):} The difference between the expected value of an estimator and the true value.

\textbf{Blue-Green Deployment:} A deployment strategy that reduces downtime by running two identical environments.

\section*{C}

\textbf{Canary Deployment:} Gradual rollout of changes to a small subset of users before full deployment.

\textbf{CI/CD:} Continuous Integration and Continuous Deployment, practices for automating software delivery.

\textbf{Concept Drift:} Changes in the underlying data distribution that affect model performance.

\textbf{Container:} A lightweight, standalone package containing code and all its dependencies.

\textbf{Cross-Validation:} A technique for assessing model generalization by partitioning data into training and validation sets.

\section*{D}

\textbf{Data Drift:} Changes in the input data distribution over time.

\textbf{Data Pipeline:} A series of data processing steps that transform raw data into a usable format.

\textbf{Data Versioning:} Tracking changes to datasets over time.

\textbf{Distributed Training:} Training models across multiple machines or GPUs.

\textbf{Docker:} A platform for developing, shipping, and running applications in containers.

\textbf{DVC (Data Version Control):} A version control system for ML projects, tracking data and models.

\section*{E}

\textbf{Embedding:} A learned representation mapping discrete variables to continuous vectors.

\textbf{Ensemble Learning:} Combining multiple models to improve performance.

\textbf{ETL (Extract, Transform, Load):} The process of extracting data from sources, transforming it, and loading it into a destination.

\textbf{Experiment Tracking:} Recording parameters, metrics, and artifacts from ML experiments.

\section*{F}

\textbf{Feature Engineering:} Creating new features from raw data to improve model performance.

\textbf{Feature Store:} A centralized repository for storing and serving features for ML models.

\textbf{Fine-Tuning:} Adapting a pre-trained model to a specific task.

\section*{G}

\textbf{GPU (Graphics Processing Unit):} Hardware accelerator commonly used for training neural networks.

\textbf{Gradient:} The vector of partial derivatives indicating the direction of steepest ascent.

\textbf{gRPC:} A high-performance RPC framework for communication between services.

\section*{H}

\textbf{Hyperparameter:} A parameter set before training that controls the learning process.

\textbf{Hyperparameter Optimization:} The process of finding optimal hyperparameter values.

\section*{I}

\textbf{IaC (Infrastructure as Code):} Managing infrastructure through code rather than manual processes.

\textbf{Inference:} Making predictions with a trained model.

\section*{K}

\textbf{Kubernetes:} An open-source container orchestration platform.

\textbf{KServe:} A Kubernetes-based platform for serving ML models.

\textbf{Kubeflow:} A machine learning toolkit for Kubernetes.

\section*{L}

\textbf{Latency:} The time taken to process a single inference request.

\textbf{Learning Rate:} A hyperparameter controlling the step size in gradient descent.

\textbf{Loss Function:} A function measuring the difference between predicted and actual values.

\section*{M}

\textbf{MLflow:} An open-source platform for the ML lifecycle.

\textbf{MLOps:} Machine Learning Operations, practices for deploying and maintaining ML systems.

\textbf{Model Registry:} A centralized repository for managing model versions.

\textbf{Model Serving:} Deploying models to make predictions in production.

\section*{N}

\textbf{Neural Architecture Search (NAS):} Automated search for optimal neural network architectures.

\section*{O}

\textbf{ONNX:} Open Neural Network Exchange, a format for representing ML models.

\textbf{Overfitting:} When a model learns training data too well, including noise, hurting generalization.

\section*{P}

\textbf{Pipeline:} A sequence of data processing or model training steps.

\textbf{Precision:} The fraction of positive predictions that are correct.

\textbf{Pruning:} Removing unnecessary weights or neurons from a neural network.

\section*{Q}

\textbf{Quantization:} Reducing model size by using lower-precision numbers.

\section*{R}

\textbf{Recall:} The fraction of actual positives correctly identified.

\textbf{Regularization:} Techniques to prevent overfitting by constraining model complexity.

\textbf{REST API:} Representational State Transfer, an architectural style for web services.

\section*{S}

\textbf{Scaling:} Adjusting feature values to a common scale.

\textbf{Shadow Deployment:} Running a new model in parallel with the existing one without affecting users.

\textbf{Sharding:} Distributing data across multiple machines.

\section*{T}

\textbf{Tensor:} A multi-dimensional array, the fundamental data structure in deep learning.

\textbf{TensorFlow Serving:} A system for serving TensorFlow models in production.

\textbf{Throughput:} The number of inference requests processed per unit time.

\textbf{TorchServe:} A tool for serving PyTorch models.

\textbf{Training-Serving Skew:} Discrepancies between training and serving environments.

\textbf{Transfer Learning:} Applying knowledge from one task to another related task.

\section*{U}

\textbf{Underfitting:} When a model is too simple to capture underlying patterns.

\section*{V}

\textbf{Validation Set:} Data used to evaluate model performance during training.

\textbf{Variance:} The amount by which model predictions vary for different training sets.

\section*{W}

\textbf{Weight Decay:} A regularization technique equivalent to L2 regularization.

\section*{Acronyms}

\begin{itemize}
    \item \textbf{AI:} Artificial Intelligence
    \item \textbf{API:} Application Programming Interface
    \item \textbf{AUC:} Area Under the Curve
    \item \textbf{AWS:} Amazon Web Services
    \item \textbf{BERT:} Bidirectional Encoder Representations from Transformers
    \item \textbf{CI/CD:} Continuous Integration/Continuous Deployment
    \item \textbf{CNN:} Convolutional Neural Network
    \item \textbf{CPU:} Central Processing Unit
    \item \textbf{CV:} Cross-Validation
    \item \textbf{DL:} Deep Learning
    \item \textbf{DVC:} Data Version Control
    \item \textbf{EDA:} Exploratory Data Analysis
    \item \textbf{ETL:} Extract, Transform, Load
    \item \textbf{GAN:} Generative Adversarial Network
    \item \textbf{GCP:} Google Cloud Platform
    \item \textbf{GPU:} Graphics Processing Unit
    \item \textbf{gRPC:} gRPC Remote Procedure Call
    \item \textbf{IaC:} Infrastructure as Code
    \item \textbf{IoU:} Intersection over Union
    \item \textbf{KNN:} K-Nearest Neighbors
    \item \textbf{LSTM:} Long Short-Term Memory
    \item \textbf{MAE:} Mean Absolute Error
    \item \textbf{MAP:} Mean Average Precision
    \item \textbf{ML:} Machine Learning
    \item \textbf{MLP:} Multi-Layer Perceptron
    \item \textbf{MLOps:} Machine Learning Operations
    \item \textbf{MSE:} Mean Squared Error
    \item \textbf{NAS:} Neural Architecture Search
    \item \textbf{NLP:} Natural Language Processing
    \item \textbf{ONNX:} Open Neural Network Exchange
    \item \textbf{REST:} Representational State Transfer
    \item \textbf{RMSE:} Root Mean Squared Error
    \item \textbf{RNN:} Recurrent Neural Network
    \item \textbf{ROC:} Receiver Operating Characteristic
    \item \textbf{SaaS:} Software as a Service
    \item \textbf{SGD:} Stochastic Gradient Descent
    \item \textbf{SLA:} Service Level Agreement
    \item \textbf{SVM:} Support Vector Machine
    \item \textbf{TPU:} Tensor Processing Unit
    \item \textbf{URL:} Uniform Resource Locator
\end{itemize}
